\begin{abstract}
% Avoid notations; avoid citations; no exceeding 250 words.
% Summarize: context, objective, methods (PRISMA), main findings, conclusion.

Equatorial Plasma Bubbles (EPBs) cause severe ionospheric scintillations
that disrupt satellite communication and navigation systems. Because
manual detection is subjective and unscalable, automated monitoring is
critical. This systematic review investigates the application of Machine
Learning (ML) techniques to automate the detection, classification, and
feature extraction of EPBs. Following the PRISMA protocol, 9 relevant
studies published between 2022 and 2025 were analyzed. The literature
reveals three methodological paradigms: end-to-end deep convolutional
neural networks, hybrid pipelines combining dimensionality reduction
with traditional classifiers, and frameworks augmented by Explainable
Artificial Intelligence. Reported detection accuracies range from 93\%
to over 99\%, with hybrid and ensemble approaches achieving the highest
performance at substantially reduced computational cost. Beyond
classification, automated morphological segmentation has enabled
solar-cycle-scale statistical profiling of EPB characteristics. The
review identifies key cross-study limitations---including environmental
contamination of optical data and training set imbalances---and
highlights critical gaps, particularly the absence of predictive
spatiotemporal models and multimodal data fusion architectures.

\end{abstract}
